% -*- coding: utf-8 -*-
\documentclass[UTF8,a4paper,10pt]{ctexart}
\usepackage[left=3.17cm, right=3.17cm, top=2.74cm, bottom=2.74cm]{geometry}
\usepackage{amsmath}
\usepackage{amssymb}
\usepackage{graphicx}
\usepackage{float}
\usepackage{caption}
\usepackage{subcaption}
\usepackage{booktabs}
\usepackage{multirow}
\usepackage{array}
\usepackage{xcolor}
\usepackage{listings}
\usepackage{algorithm}
\usepackage{algorithmicx}
\usepackage{algpseudocode}
\usepackage[colorlinks,linkcolor=black,anchorcolor=blue,urlcolor=blue]{hyperref}
\graphicspath{{../figures/}{figures/}}

\ctexset{
  section = {name = {,、}, number = \chinese{section}},
  subsection = {name = {（,）}, number = \chinese{subsection}},
  subsubsection = {name = {,.}, number = \arabic{subsubsection}}
}

\usepackage{fancyhdr}
\pagestyle{fancy}
\lhead{\kaishu \leftmark}
\rhead{\kaishu 高级计算机系统大作业报告}
\cfoot{\thepage}
\renewcommand{\headrulewidth}{0.1pt}
\renewcommand{\footrulewidth}{0pt}

\floatname{algorithm}{Algorithm}
\renewcommand{\algorithmicrequire}{\textbf{Input:}}
\renewcommand{\algorithmicensure}{\textbf{Output:}}

\lstset{
  language=bash,
  basicstyle=\ttfamily\small,
  keywordstyle=\color{blue!70}\bfseries,
  commentstyle=\color{green!40!black},
  stringstyle=\color{orange!70!black},
  numbers=left,
  numberstyle=\tiny\color{gray},
  frame=single,
  breaklines=true,
  columns=fullflexible,
  tabsize=4
}

\newcommand{\HRule}{\rule{\linewidth}{0.5mm}}
\newcommand{\TTFT}{\mathrm{TTFT}}
\newcommand{\Saved}{\Delta T}
\newcommand{\FigurePlaceholder}[1]{%
  \fbox{\begin{minipage}[c][0.26\textheight][c]{0.88\textwidth}
    \centering
    图表占位：\texttt{\detokenize{#1}}
  \end{minipage}}%
}
\newcommand{\MaybeFigure}[2][]{%
  \IfFileExists{#2}{%
    \includegraphics[#1]{#2}%
  }{%
    \IfFileExists{../figures/#2}{%
      \includegraphics[#1]{#2}%
    }{%
      \IfFileExists{figures/#2}{%
        \includegraphics[#1]{#2}%
      }{%
        \FigurePlaceholder{#2}%
      }%
    }%
  }%
}

\begin{document}

\begin{titlepage}
    \begin{center}
    \IfFileExists{NKU.png}{\includegraphics[width=0.78\textwidth]{NKU.png}\\[1cm]}{}
    \textsc{\Huge \kaishu{\textbf{南\ \ \ \ \ \ 开\ \ \ \ \ \ 大\ \ \ \ \ \ 学}} }\\[0.9cm]
    \textsc{\huge \kaishu{\textbf{计\ \ 算\ \ 机\ \ 学\ \ 院}}}\\[0.5cm]
    \textsc{\Large \textbf{高级计算机系统大作业报告}}\\[0.8cm]
    \HRule \\[0.9cm]
    { \LARGE \bfseries 大模型推理前缀缓存与在线 Offload 策略实验}\\[0.4cm]
    \HRule \\[2.0cm]
    \centering
    \textsc{\LARGE 丁益三 2310724\kaishu{\ \ \ \ }}\\[0.5cm]
    \textsc{\LARGE \kaishu{年级\ :\ 2023级}}\\[0.5cm]
    \textsc{\LARGE \kaishu{专业\ :\ 计算机科学与技术}}\\[0.5cm]
    \textsc{\LARGE \kaishu{指导教师\ :\ 王刚}}\\[0.5cm]
    \vfill
    {\Large \today}
    \end{center}
\end{titlepage}

\newpage
\thispagestyle{empty}
\renewcommand{\abstractname}{\kaishu \large \textbf{摘要}}
\begin{abstract}
本次大作业围绕大模型推理中的 Prefix Cache 和 KV Cache offload 展开。我的实现没有从一开始就直接做 offload 策略，而是先建立可复现的 vLLM 服务模式 baseline，再补充请求级 cache event 观测，最后在真实 CoQA trace 上做严格 online 的 CARGO 策略 replay。这样的顺序主要是为了避免两个常见误判：第一，把每次请求重新启动模型进程得到的延迟误认为服务场景下的 prefix-cache 行为；第二，在策略 replay 中提前使用 future reuse 信息，使策略看起来有效但无法真实部署。

实验结果显示，vLLM 原生 Prefix Cache 在长前缀、高复用场景下确实能显著降低首 token 延迟。以 CoQA 1536-token 前缀为例，warm 请求相对关闭 Prefix Cache 的 median TTFT saved 为 38.37ms，speedup 约为 2.08x；但当 prefix group 数从 16 增至 24 后，GPU-only cache 的 revisit useful-hit rate 从约 1.0 跌至 0，hit coverage 下降到约 0.021，同时出现数千次 eviction。这说明 GPU 内部 prefix cache 的问题不是“能否命中一个 block”，而是在容量压力下能否保住足够长的有用前缀。

基于这一现象，我把 offload 策略抽象为两级缓存 replay。CARGO 策略只使用当前请求之前的访问历史、prefix 长度、估计 KV 大小、最近访问时间和固定 first-access prior，不读取 \texttt{planned\_reuse\_count}、\texttt{phase} 或未来访问位置。策略的核心是用“单位字节的期望 TTFT 净收益”决定 GPU/CPU placement，并用 coverage guard 过滤只有短 header 命中的假阳性。在 CoQA popularity-skew trace 上，online CARGO 在 g24+g32+g48 设置下恢复 TTFT saved 15.20s，优于 LRU/LFU 的 14.24s 和 size-aware 的 12.93s；在 mixed prefix length trace 上，CARGO 恢复 6.93s，而 best conventional 只有 2.86s。这个差异说明，offload 策略不能只追求 hit 数量，还要同时考虑 prefix 的长度、reload 成本和一次命中能节省多少 prefill。

我也尝试了 LMCache MP mode 的真实 CPU offload baseline。该路径在隔离环境中成功跑通，LMCache server 能观察到跨进程 L1 prefix hit；但在 WSL2 上 CUDA IPC 出现 \texttt{invalid resource handle}，只能退回 non-GPU transfer，导致 g24/g32 revisit TTFT 仍在 338--350ms 左右。因此，本报告没有把 LMCache 结果包装成完整性能收益，而把它作为真实系统部署风险和后续 native Linux 实验的依据。

\par\noindent\textbf{关键词：} 大模型推理；vLLM；Prefix Cache；KV Cache；CoQA；在线策略；CARGO；Offload
\end{abstract}

\tableofcontents
\newpage
\setcounter{page}{1}

\textbf{项目仓库：} \texttt{https://github.com/Ding-1203/kv-cache-offload}

\section{实验背景与任务理解}

大模型推理通常可以分为 prefill 和 decode 两个阶段。prefill 阶段需要对 prompt 中的 token 逐层计算 Key/Value 张量；decode 阶段则在已有 KV Cache 的基础上逐 token 生成。对于长上下文问答、文档检索增强、代码补全等应用，多个请求往往共享相同的 system prompt、文档正文或 story，只在最后的问题部分不同。如果每个请求都重复计算共享前缀的 KV，prefill 会成为明显浪费；Prefix Cache 的作用就是把已经计算过的共享前缀 KV 保存下来，在后续请求中直接复用。

本题目要求关注的是“前缀缓存与换入换出策略”，而不是单纯调用一个现成开关。我的理解是，作业至少包含三层问题。第一层是系统问题：在真实推理服务中，Prefix Cache 是否真的跨请求生效，TTFT 是否下降。第二层是观测问题：当 latency 改变时，内部 cache hit、hit tokens 和 eviction 是否能解释这个变化。第三层是策略问题：当 GPU KV cache 容量不足时，哪些前缀值得继续放在 GPU，哪些可以 offload 到 CPU，哪些即使 offload 也不值得 reload。

因此，我的实际工作按“先基线、再观测、后策略”的顺序推进。前半部分使用 vLLM OpenAI-compatible server，模型只启动一次并常驻 GPU，runner 连续发送请求，以保证 Prefix Cache 可以跨请求复用。后半部分在 CoQA trace 上构造容量压力，并用 replay 的方式验证 CARGO 这样的 utility-aware 策略是否优于 LRU、LFU 和 size-aware 等常规策略。这里刻意把 replay 设为 online：除了 oracle upper bound 外，所有策略只能看到当前请求之前的信息。这样虽然会降低策略上限，但更接近真实系统中可以部署的决策条件。

\section{实验环境与总体框架}

\subsection{软硬件环境}

实验在 WSL2 Ubuntu 22.04 环境中完成，核心软件版本和硬件条件如表 \ref{tab:env} 所示。所有 vLLM 实验均优先使用本地模型路径 \texttt{QWEN05B\_PATH}，避免重复下载模型。由于显存只有 8GB，我选择 Qwen2.5-0.5B-Instruct 作为主模型。这个模型足够小，可以在本地稳定运行；同时它仍然具有真实 transformer KV Cache 结构，能够反映 prefix reuse 和 capacity pressure 对 TTFT 的影响。

\begin{table}[H]
\centering
\small
\caption{实验环境配置}
\label{tab:env}
\begin{tabular}{ll}
\toprule
项目 & 配置 \\
\midrule
操作系统 & WSL2 Ubuntu 22.04 \\
GPU & NVIDIA RTX 4060 Laptop GPU, 8GB \\
模型 & Qwen2.5-0.5B-Instruct，本地路径由 \texttt{QWEN05B\_PATH} 提供 \\
Python & 3.12.13 \\
vLLM & 0.20.0+cu129 \\
PyTorch & 2.11.0+cu129 \\
Triton & 3.6.0 \\
主要数据 & CoQA train/validation 派生的 story-question 请求 \\
\bottomrule
\end{tabular}
\end{table}

\subsection{实验框架}

整体框架由 workload builder、runner、observer 和 analyzer 四部分组成。workload builder 将 CoQA story 和 question 转成 JSONL 请求，每行包含 \texttt{request\_id}、\texttt{prefix\_group}、prompt、prefix 长度和场景标签。runner 逐行读取 JSONL，通过 vLLM OpenAI API 发送请求，并记录 TTFT、total latency、prompt tokens、开始时间和结束时间。observer 通过 \texttt{sitecustomize.py} 和 monkey patch 注入 vLLM 进程，记录 prefix lookup、prefix hit/miss、block touch、store、free 和 evict 事件。analyzer 最后把 latency CSV 与 cache event JSONL 对齐，生成 summary 和图表。

这个框架里最关键的实验约束是服务模式。vLLM server 在每组实验中只启动一次，模型常驻 GPU，runner 在同一服务进程生命周期内连续发送多个请求。只有这样，上一条请求写入的 KV block 才可能被下一条请求复用。如果把每个请求写成独立 Python 进程并反复初始化模型，那么 Prefix Cache 每次都会被清空，实验结论会偏离真实在线服务。

\begin{algorithm}[H]
\caption{服务模式下的 Prefix Cache 实验 runner}
\label{alg:runner}
\begin{algorithmic}[1]
\Require JSONL workload, vLLM OpenAI endpoint, model name
\Ensure 每个请求的 latency CSV
\State 启动一次 vLLM server，使模型和 Prefix Cache 常驻服务进程
\For{workload 中的每一行请求 $r$}
    \State 读取 prompt、prefix group、request id 和实验标签
    \State 通过 OpenAI-compatible API 发送请求，并在 extra body 中传入 request id
    \State 记录请求开始时间、首 token 到达时间、结束时间和状态
    \State 将 TTFT、total latency、prompt tokens、prefix group、scenario 写入 CSV
\EndFor
\State 如开启 observer，则按 request id 和时间窗口将 cache event 与 latency CSV 对齐
\end{algorithmic}
\end{algorithm}

\section{指标定义与问题建模}

\subsection{TTFT 与 Prefix Cache 收益}

本报告主要使用 TTFT（time to first token）作为延迟指标。对于交互式问答服务，用户最直接感受到的是首 token 出现时间，而 Prefix Cache 主要减少的是 prefill 计算，因此 TTFT 比 total latency 更能体现前缀复用收益。对一个请求 $q$，可以把 TTFT 粗略分解为
\begin{equation}
\TTFT(q) \approx T_{\mathrm{sched}}(q) + T_{\mathrm{lookup}}(q)
    + T_{\mathrm{prefill}}\!\left(L_q-H_q\right)
    + T_{\mathrm{decode},1}(q) + \epsilon_q .
\label{eq:ttft}
\end{equation}
其中 $L_q$ 是 prompt 中需要 prefill 的前缀 token 数，$H_q$ 是被 Prefix Cache 命中的 token 数。Prefix Cache 的理想作用是增大 $H_q$，从而减少实际 prefill 的 token 数 $L_q-H_q$。由于调度、网络、本地进程负载等因素会带来噪声，单个请求的 TTFT 并不稳定，所以后续表格主要报告 median 或按场景聚合后的结果。

我将开启和关闭 Prefix Cache 的 TTFT 差定义为
\begin{equation}
\Saved(q) = \TTFT_{\mathrm{off}}(q) - \TTFT_{\mathrm{on}}(q).
\label{eq:saved}
\end{equation}
当 $\Saved(q)>0$ 时，开启 Prefix Cache 降低了首 token 延迟；当 $\Saved(q)\le 0$ 时，说明该请求没有从缓存复用中得到可观收益，或者收益被系统噪声抵消。这个定义也用于后续 replay 中估计某个 prefix 被命中时可能节省的 prefill 时间。

\subsection{Hit coverage 与 useful hit}

只记录是否发生 prefix hit 并不够。vLLM 的 prefix cache 以 block 为单位工作，有时请求可能命中很短的公共 header，但真正的大段 story 前缀已经被 evict。此时 binary hit 为真，但 TTFT 仍接近 cold miss。为避免这种假阳性，我使用 hit coverage 衡量命中 token 占目标前缀长度的比例：
\begin{equation}
\rho(q) = \frac{H_q}{L^p_q},
\qquad
\mathrm{useful}(q)=\mathbf{1}\left[H_q \ge 32 \land \rho(q)\ge \theta\right].
\label{eq:coverage}
\end{equation}
其中 $L^p_q$ 是该请求所属 prefix group 的目标前缀长度，$\theta$ 在实验中取 $0.5$。额外的 32-token 下界用于过滤 system prompt 或短 header 带来的伪命中。下文所有 useful-hit rate 均按公式 \ref{eq:coverage} 计算；后续章节只引用该定义，用它区分“发生过短命中”和“命中了足够长的可复用前缀”。

\subsection{KV Cache 大小估计}

Offload 策略需要比较“保留某个 prefix 的收益”和“它占用的空间”。我在 replay 中以 prefix group 为粒度估计 KV 大小，而不是模拟 vLLM 内部每个 block 的张量布局。对长度为 $L_p$ 的 prefix，估计字节数为
\begin{equation}
B(p)=
\left\lceil \frac{L_p}{b} \right\rceil b
\cdot 2 \cdot N_{\mathrm{layer}} \cdot N_{\mathrm{kv\_head}}
\cdot d_{\mathrm{head}} \cdot s_{\mathrm{dtype}} .
\label{eq:kvbytes}
\end{equation}
其中 $b$ 是 vLLM block size，系数 2 对应 Key 和 Value。对本实验中的 Qwen2.5-0.5B-Instruct，近似取 $b=16$，$N_{\mathrm{layer}}=24$，$N_{\mathrm{kv\_head}}=2$，$d_{\mathrm{head}}=64$，FP16/BF16 的 $s_{\mathrm{dtype}}=2$ bytes。因此每个 token 的 KV 大小约为
\[
2 \times 24 \times 2 \times 64 \times 2 = 12288\ \mathrm{bytes},
\]
1536-token prefix 约占 $18.0$ MiB。这个数量级解释了为什么在 8GB 笔记本 GPU 上，prefix group 数稍微增加就会触发明显的 cache eviction。

\section{vLLM Prefix Cache Baseline}

\subsection{从构造 workload 到 CoQA workload}

第一阶段我先用构造 workload 验证实验框架是否能稳定复现 Prefix Cache 收益，包括 Prefix Cache ON/OFF、不同 prefix 长度、不同复用率和多 prefix group 轮换访问。这一步的价值在于排除工程层错误，例如服务没有复用、请求格式不一致、max model length 设置过小等。随后我将主要实验换成 CoQA。CoQA 每条样本包含一个 story 和围绕该 story 的多轮 question，天然对应“相同长前缀 + 不同问题”的请求模式，比随机拼接文本更接近真实问答服务。

CoQA workload 分为三类。第一类是 prefix-length sweep，固定复用结构但改变 story 前缀长度，用来观察长前缀是否带来更大 TTFT saved。第二类是 popularity-skew，不同 story 的访问热度不同，用来观察热前缀是否贡献了更多累计收益。第三类是 capacity-pressure，先 seed 若干 prefix group，再插入大量 filler group，最后 revisit seed group，用来观察 GPU-only Prefix Cache 在容量压力下是否发生 cliff。

\subsection{Prefix 长度决定可节省的 prefill}

表 \ref{tab:prefix-length} 展示了 CoQA prefix-length sweep 中 warm 请求的 median TTFT。这里的 \texttt{Prefix ON} 表示 vLLM Prefix Cache 开启时的 TTFT，\texttt{Prefix OFF} 表示关闭 Prefix Cache 的对照结果。512-token 档位没有稳定收益，median saved 为 $-1.82$ms，这主要说明短前缀的可节省 prefill 不足以稳定超过调度和测量噪声。从 768 token 开始，saved TTFT 随 prefix 长度增加而持续扩大；到 1536 token 时，median saved 为 38.37ms，speedup 约为 2.08x。

\begin{table}[H]
\centering
\small
\caption{CoQA prefix length sweep：warm 请求的 TTFT 对比}
\label{tab:prefix-length}
\begin{tabular}{cccccc}
\toprule
Prefix bucket & 请求数 & Prefix ON TTFT(ms) & Prefix OFF TTFT(ms) & Saved(ms) & Speedup \\
\midrule
512  & 120 & 35.73 & 33.72 & -1.82 & 0.95x \\
768  & 120 & 33.81 & 43.75 & 10.03 & 1.30x \\
1024 & 120 & 33.44 & 54.96 & 21.35 & 1.63x \\
1280 & 120 & 34.33 & 65.15 & 30.24 & 1.89x \\
1536 & 120 & 36.09 & 74.93 & 38.37 & 2.08x \\
\bottomrule
\end{tabular}
\end{table}

图 \ref{fig:coqa-prefix-length} 对这一趋势做了可视化。更重要的不是单个点的数值，而是图中左右轴的语义关系：当 saved 为 0ms 时，speedup 应接近 1.0x；当 prefix 更长时，关闭 Prefix Cache 的 prefill 成本继续增加，而开启 Prefix Cache 后 warm 请求的 TTFT 基本维持在 33--36ms 区间。这个形态说明 Prefix Cache 的收益来自减少重复 prefill，而不是 decode 或网络开销的偶然变化。

\begin{figure}[H]
\centering
\MaybeFigure[width=0.92\textwidth]{cache_observation/coqa_final/01_coqa_prefix_length_ttft_savings.png}
\caption{CoQA prefix length sweep：prefix 越长，Prefix Cache 对 TTFT 的节省越明显。该图对应真实输出 \texttt{figures/cache\_observation/coqa\_final/01\_coqa\_prefix\_length\_ttft\_savings.png}。}
\label{fig:coqa-prefix-length}
\end{figure}

\subsection{Cold 与 warm 请求必须分开看}

表 \ref{tab:on-off-state} 将 CoQA prefix-length 和 popularity-skew 两类实验按 cold/warm 分开统计。cold 请求是某个 prefix group 的首次访问，此时没有可复用 KV；warm 请求是同一 prefix group 的后续访问。结果显示，cold 请求几乎没有收益，甚至会因为系统噪声出现轻微负收益；warm 请求则稳定降低 TTFT。在 popularity-skew 中，warm saved 达到 39.44ms，和 1536-token 长前缀的结果一致。

\begin{table}[H]
\centering
\small
\caption{CoQA 实验中 cold/warm 请求的 Prefix Cache ON/OFF 对比}
\label{tab:on-off-state}
\begin{tabular}{llccccc}
\toprule
实验类型 & 状态 & 请求数 & ON TTFT(ms) & OFF TTFT(ms) & Saved(ms) & Useful hit \\
\midrule
Prefix length & cold & 200 & 55.78 & 54.65 & -0.04 & 0.00 \\
Prefix length & warm & 600 & 34.53 & 54.97 & 20.82 & 1.00 \\
Popularity skew & cold & 92 & 74.92 & 75.13 & -0.09 & 0.00 \\
Popularity skew & warm & 272 & 36.08 & 75.48 & 39.44 & 1.00 \\
\bottomrule
\end{tabular}
\end{table}

图 \ref{fig:coqa-on-off} 的作用是给 baseline 结论加上边界条件。Prefix Cache 不是一个无条件的加速开关；它的收益只在“相同前缀已经被计算过，并且后续请求仍能命中足够长的缓存前缀”时出现。这个结论也直接引出下一阶段的问题：如果 GPU cache 容量不足，warm 请求会重新退化成 cold-like 请求。

\begin{figure}[H]
\centering
\MaybeFigure[width=0.92\textwidth]{cache_observation/coqa_final/02_coqa_on_off_by_cache_state.png}
\caption{CoQA cold/warm 请求下 Prefix Cache ON/OFF 的 TTFT 对比。该图对应真实输出 \texttt{figures/cache\_observation/coqa\_final/02\_coqa\_on\_off\_by\_cache\_state.png}。}
\label{fig:coqa-on-off}
\end{figure}

\section{Prefix Cache 观测与容量压力}

\subsection{vLLM 内部事件观测}

为了避免只从外部 latency 推断缓存行为，我实现了一个轻量级 observer。服务启动脚本将 \texttt{PYTHONPATH} 指向项目的 \texttt{src/}，并在设置 \texttt{VLLM\_CACHE\_OBSERVER=1} 后让 vLLM API 进程和 EngineCore 子进程在导入 \texttt{sitecustomize.py} 时自动调用 \texttt{vllm\_cache\_observer.install()}。这个 observer 只包裹 cache 管理层的方法，不修改调度、分配或 eviction 策略，因此不会改变被观测系统本身的行为。

具体插桩分成两层。第一层在 \texttt{KVCacheManager.get\_computed\_blocks} 外包装请求级 prefix lookup，记录 \texttt{prefix\_lookup}，并根据 \texttt{hit\_tokens} 进一步写入 \texttt{prefix\_hit} 或 \texttt{prefix\_miss}；同时保存 vLLM 内部 request id、归一化后的 workload request id、hit tokens、hit blocks、prompt tokens 以及命中的 block hash。\texttt{allocate\_slots} 和 \texttt{free} 主要用于维护当前请求上下文，使后续 block 级事件能够带上对应的 request id。第二层在 \texttt{BlockPool.cache\_full\_blocks}、\texttt{touch}、\texttt{free\_blocks} 和 \texttt{\_maybe\_evict\_cached\_block} 上记录 \texttt{block\_store}、\texttt{block\_touch}、\texttt{block\_free} 和 \texttt{block\_evict}，字段包括 block id、block hash 和 ref count。

这些事件统一由 JSONL logger 追加写入，基础字段包括 \texttt{timestamp}、\texttt{pid}、\texttt{event\_type}、\texttt{internal\_request\_id}、\texttt{request\_id}、\texttt{hit\_tokens}、\texttt{hit\_blocks}、\texttt{block\_id}、\texttt{block\_hash} 和 \texttt{ref\_cnt}。前文已经用公式 \ref{eq:coverage} 定义了 hit coverage 和 useful-hit 判定；本节的 observer 主要提供计算这些指标所需的 $H_q$、block 事件和 eviction 证据，并通过 request id 与 runner CSV 对齐，从而解释 TTFT 曲线背后的真实 cache 行为。

\subsection{热度与累计收益的关系}

Popularity-skew workload 用不同 planned reuse count 的 CoQA story 模拟热度差异。表 \ref{tab:popularity} 展示了热度越高，warm saved 和 useful-hit rate 越稳定；当 planned reuse count 从 1 增加到 8 或 12 时，useful-hit rate 从 0.50 提升到 0.89/0.92，warm saved 也维持在约 39ms。这里 planned reuse count 只用于 workload 构造和结果分组，后续 online 策略不会直接读取它。

\begin{table}[H]
\centering
\small
\caption{CoQA popularity-skew：访问热度与 warm 收益}
\label{tab:popularity}
\begin{tabular}{ccccc}
\toprule
Planned reuse count & Prefix groups & 请求数 & Warm saved(ms) & Useful-hit rate \\
\midrule
1  & 32 & 64 & 5.65  & 0.50 \\
2  & 32 & 96 & 36.26 & 0.67 \\
4  & 16 & 80 & 38.67 & 0.80 \\
8  & 8  & 72 & 39.84 & 0.89 \\
12 & 4  & 52 & 39.51 & 0.92 \\
\bottomrule
\end{tabular}
\end{table}

图 \ref{fig:coqa-pareto} 进一步说明，prefix cache 的收益具有明显集中性。少量高复用 prefix group 贡献了较多累计 saved TTFT。这正是 utility-aware offload 的动机：如果 GPU 放不下所有 KV，保留“更可能再次访问、且一次命中能节省更多 prefill”的 prefix，比机械地保留最近访问或最小对象更合理。

\begin{figure}[H]
\centering
\MaybeFigure[width=0.92\textwidth]{cache_observation/coqa_final/03_coqa_popularity_utility_pareto.png}
\caption{CoQA popularity-skew：prefix 热度、累计收益和 utility 集中性。该图对应真实输出 \texttt{figures/cache\_observation/coqa\_final/03\_coqa\_popularity\_utility\_pareto.png}。}
\label{fig:coqa-pareto}
\end{figure}

\subsection{容量 pressure cliff}

容量压力实验采用 seed--filler--revisit 访问模式。先访问 $g$ 个 seed prefix group，使它们进入 cache；随后访问 $4g$ 个 filler group，占用大量 KV 空间；最后再次访问 seed group，观察 seed 是否仍能命中。这个设计刻意模拟真实服务中的“热点文档被其他请求冲掉后再次访问”的场景。

表 \ref{tab:pressure} 是最关键的 baseline 结果。g8 和 g16 下，revisit 请求仍然基本保持 useful hit，median TTFT 在 39--42ms 左右；但从 g24 开始，revisit TTFT 突然上升到约 76ms，useful-hit rate 变为 0，hit coverage 只有 0.0208，同时观测到数千次 eviction。这个 cliff 说明，GPU-only Prefix Cache 在容量足够时表现很好，但一旦超过某个工作集阈值，就会把原本 warm 的请求重新变成 cold-like 请求。

\begin{table}[H]
\centering
\small
\caption{CoQA capacity-pressure：revisit 阶段的 cache cliff}
\label{tab:pressure}
\begin{tabular}{cccccc}
\toprule
Prefix groups & Revisit TTFT(ms) & P90 TTFT(ms) & Useful-hit rate & Hit coverage & Evictions \\
\midrule
8  & 41.63 & 44.05 & 1.00 & 1.0000 & 0 \\
16 & 39.36 & 47.40 & 1.00 & 0.9993 & 0 \\
24 & 75.94 & 79.72 & 0.00 & 0.0208 & 3496 \\
32 & 76.27 & 79.09 & 0.00 & 0.0208 & 3033 \\
48 & 76.00 & 78.89 & 0.00 & 0.0208 & 4641 \\
\bottomrule
\end{tabular}
\end{table}

图 \ref{fig:coqa-pressure} 将这个 cliff 同时画在 TTFT、hit coverage 和 eviction 维度上。相比单独报告 TTFT，这张图更清楚地说明了性能退化的原因：不是模型突然变慢，而是 prefix cache 中可复用的长前缀被容量压力驱逐了。由此可以得到后续 offload 的核心问题：如果 GPU 放不下所有 prefix，是否可以把被驱逐但未来可能有价值的 prefix 放到 CPU，并在 revisit 时用低于 cold prefill 的成本 reload。

\begin{figure}[H]
\centering
\MaybeFigure[width=0.92\textwidth]{cache_observation/coqa_final/04_coqa_pressure_cliff.png}
\caption{CoQA capacity-pressure：g16 到 g24 之间出现明显 cache cliff，revisit TTFT 上升且 useful-hit rate 归零。该图对应真实输出 \texttt{figures/cache\_observation/coqa\_final/04\_coqa\_pressure\_cliff.png}。}
\label{fig:coqa-pressure}
\end{figure}

\section{Online CARGO Replay 策略}

\subsection{为什么不能直接用未来信息}

在最初的 replay 设计中，一个容易犯的错误是使用 workload 中的 planned reuse count 或 phase 标签来指导策略。这样的策略在离线分析中可能表现很好，但真实系统在当前时刻并不知道某个 prefix 之后还会被访问几次，也不知道当前请求属于 seed、filler 还是 revisit。因此，我在最终版本中把策略分为两类。非 oracle 策略只能使用历史信息，包括当前请求之前的访问次数、最近访问时间、prefix 长度、估计 KV 大小和历史观测得到的 saved curve。Oracle upper bound 可以使用 future reuse，但它只作为上界，不作为可部署策略。

这个约束使实验结论更保守，也更有意义。如果 online CARGO 仍然优于 LRU、LFU 和 size-aware，说明它的优势来自对历史信号、对象大小和 reload 成本的建模，而不是偷看未来。

\subsection{策略建模}

Replay 中的缓存粒度是 prefix group。GPU 层表示 vLLM 内部可快速命中的 KV；CPU 层表示被 offload 后仍可 reload 的 KV。由于我没有修改 vLLM 内部 KV tensor reload，本阶段 replay 不模拟 block-level tensor 传输，而是用一个成本模型估计 CPU reload 的收益。对 prefix $p$，从 baseline 得到的长度相关 saved curve 记为 $S(L_p)$，CPU reload 成本估计为
\begin{equation}
C_{\mathrm{cpu}}(p) = c_0 + \frac{B(p)}{\mathrm{BW}_{\mathrm{cpu}}},
\label{eq:reload-cost}
\end{equation}
其中 $c_0$ 是固定调度开销，实验中取 8ms，$\mathrm{BW}_{\mathrm{cpu}}$ 取 2GB/s。则一次 CPU reload 的净收益为
\begin{equation}
G_{\mathrm{reload}}(p) =
\max\left(0, S(L_p) - C_{\mathrm{cpu}}(p)\right).
\label{eq:net-gain}
\end{equation}
如果 reload 成本已经超过命中该 prefix 能节省的 prefill 时间，那么即使 CPU 中保存了该 prefix，也不应该为了它付出 reload 开销。

CARGO 的核心打分是单位字节 utility。令 $n_p(t-1)$ 为当前时刻前 prefix $p$ 已经被访问的次数，历史热度信号为
\begin{equation}
A_p(t)=\alpha+\max(0,n_p(t-1)-1),
\label{eq:history-signal}
\end{equation}
其中 $\alpha=0.25$ 是 first-access prior，用来避免新 prefix 完全没有进入 CPU 的机会。GPU 侧保留价值和 CPU 侧 reload/offload 价值分别定义为
\begin{equation}
U_{\mathrm{gpu}}(p,t)=\frac{S(L_p)\cdot A_p(t)}{B(p)},
\qquad
U_{\mathrm{cpu}}(p,t)=\frac{G_{\mathrm{reload}}(p)\cdot A_p(t)}{B(p)}.
\label{eq:cargo-score}
\end{equation}
当 GPU 容量不足时，CARGO 优先驱逐 $U_{\mathrm{gpu}}$ 较低的 prefix；被驱逐对象只有在 $U_{\mathrm{cpu}}>0$ 时才进入 CPU 层。这样做的直觉是，长 prefix 一次命中可能节省更多 prefill，但它也占用更多字节；高热度 prefix 更值得保留，但如果 reload 成本过高，放到 CPU 也未必划算。

\begin{algorithm}[H]
\caption{Online CARGO replay 决策流程}
\label{alg:cargo}
\begin{algorithmic}[1]
\Require 请求序列 $q_1,\ldots,q_T$，GPU 容量 $C_g$，CPU 容量 $C_c$
\Ensure 每个策略的 recovered TTFT saved、GPU hit、CPU reload 和 cold miss
\State 初始化 GPU cache、CPU cache 和每个 prefix 的历史计数
\For{$t=1$ to $T$}
    \State 令 $p$ 为请求 $q_t$ 的 prefix group
    \If{$p$ 在 GPU cache 且满足 coverage guard}
        \State 记录 GPU hit，获得 $S(L_p)$ 的 saved
    \ElsIf{$p$ 在 CPU cache 且 $G_{\mathrm{reload}}(p)>0$}
        \State 记录 CPU reload，获得 $G_{\mathrm{reload}}(p)$ 的净 saved
        \State 将 $p$ 提升到 GPU cache
    \Else
        \State 记录 cold miss
    \EndIf
    \While{GPU cache 占用超过 $C_g$}
        \State 按 $U_{\mathrm{gpu}}$ 选择最低价值 prefix $v$ 驱逐
        \If{$U_{\mathrm{cpu}}(v)>0$ 且 CPU cache 有可用价值}
            \State 将 $v$ 放入 CPU cache；必要时驱逐 CPU 中最低 $U_{\mathrm{cpu}}$ 的 prefix
        \EndIf
    \EndWhile
    \State 使用请求 $q_t$ 更新 $p$ 的访问次数和最近访问时间
\EndFor
\end{algorithmic}
\end{algorithm}

对照策略包括 GPU-only、CPU LRU、CPU LFU、size-aware、CARGO 和 oracle upper bound。LRU 只看最近访问时间，适合时间局部性强但对象大小相近的场景；LFU 只看历史次数，可能保留很久以前的热对象；size-aware 倾向于保留小对象，但可能把 saved 很低的小 prefix 留下来；CARGO 则试图同时考虑 saved、bytes、历史热度和 reload 成本。Oracle 使用 future reuse，因此只代表上限。

\subsection{Replay 结果分析}

表 \ref{tab:cargo-summary} 汇总了 g24、g32 和 g48 三个超过 GPU-only cliff 的容量压力场景。Popularity-skew trace 中，GPU-only 只能恢复 9.06s saved，因为很多 revisit 已经变成 cold miss；LRU/LFU 通过 CPU 层恢复到 14.24s；CARGO 进一步提升到 15.20s。这个提升并不夸张，因为 popularity-skew 中 prefix 长度较一致，LRU/LFU 已经能利用访问热度。CARGO 的优势主要来自减少低价值对象占位，并在 CPU reload 前扣除传输成本。

Mixed-length trace 更能体现 utility-aware 的必要性。LRU/LFU/GPU-only 在该 trace 上几乎无法恢复 saved；size-aware 的 useful-hit rate 反而达到 0.380，但 recovered saved 只有 2.86s。这说明“命中次数多”并不等价于“节省时间多”，因为它可能命中的是短 prefix 或低 saved prefix。CARGO 的 useful-hit rate 只有 0.197，却恢复了 6.93s saved，说明它更倾向于选择少量高价值 prefix，而不是追求表面 hit rate。

\begin{table}[H]
\centering
\small
\caption{Online replay 在 g24+g32+g48 下的策略对比}
\label{tab:cargo-summary}
\begin{tabular}{llccccc}
\toprule
Trace & 策略 & Recovered saved(s) & Useful rate & GPU hit & CPU reload & Cold miss \\
\midrule
\multirow{5}{*}{Popularity skew}
& GPU-only & 9.06 & 0.289 & 236 & 0   & 580 \\
& LRU/LFU  & 14.24 & 0.593 & 236 & 248 & 332 \\
& Size-aware & 12.93 & 0.578 & 176 & 296 & 344 \\
& CARGO & 15.20 & 0.612 & 273 & 226 & 317 \\
& Oracle upper bound & 24.61 & 0.885 & 545 & 177 & 94 \\
\midrule
\multirow{5}{*}{Mixed length}
& GPU-only & 0.00 & 0.000 & 0 & 0   & 1800 \\
& LRU/LFU  & 0.00 & 0.000 & 0 & 0   & 1800 \\
& Size-aware & 2.86 & 0.380 & 576 & 108 & 1116 \\
& CARGO & 6.93 & 0.197 & 25 & 329 & 1446 \\
& Oracle upper bound & 16.56 & 0.374 & 300 & 374 & 1126 \\
\bottomrule
\end{tabular}
\end{table}

图 \ref{fig:no-future} 是 online 约束的 sanity check。它用于确认非 oracle 策略没有读取未来字段，并展示 CARGO 与 oracle 之间仍然存在明显差距。如果 CARGO 接近 oracle，需要怀疑是否泄漏了 future reuse；而当前结果中 CARGO 明显低于 oracle，但稳定高于 conventional baseline，这更符合一个可部署在线策略的预期。

\begin{figure}[H]
\centering
\MaybeFigure[width=0.88\textwidth]{offload_online/01_no_future_sanity.png}
\caption{Online replay 的 no-future sanity check。该图对应真实输出 \texttt{figures/offload\_online/01\_no\_future\_sanity.png}。}
\label{fig:no-future}
\end{figure}

图 \ref{fig:policy-gain} 和图 \ref{fig:capacity-sweep} 分别从策略收益和容量变化角度展示结果。前者强调 CARGO 在两个 trace 上都不依赖 future 信息仍能超过常规策略；后者显示随着压力增大，GPU-only 的 saved 受限，而 CARGO 可以通过 CPU 层恢复一部分 cliff 后损失的收益。不过 CARGO 和 oracle 之间仍有差距，说明单纯依靠历史计数和静态 saved curve 还不足以完全预测未来访问，后续可以进一步加入时间衰减、队列状态和应用侧会话信息。

\begin{figure}[H]
\centering
\MaybeFigure[width=0.88\textwidth]{offload_online/02_policy_gain.png}
\caption{Online CARGO 与 LRU、LFU、size-aware、GPU-only 和 oracle 的 recovered TTFT saved 对比。该图对应真实输出 \texttt{figures/offload\_online/02\_policy\_gain.png}。}
\label{fig:policy-gain}
\end{figure}

\begin{figure}[H]
\centering
\MaybeFigure[width=0.88\textwidth]{offload_online/03_capacity_sweep.png}
\caption{不同 pressure level 下的策略收益变化。该图对应真实输出 \texttt{figures/offload\_online/03\_capacity\_sweep.png}。}
\label{fig:capacity-sweep}
\end{figure}

图 \ref{fig:decision-audit} 用于解释策略行为本身。它不是只报告最终 saved，而是把 GPU hit、CPU reload 和 cold miss 拆开。这个拆分能看出 CARGO 的策略特征：它没有盲目增加 CPU reload 次数，而是在 reload 成本模型允许的情况下选择更高净收益的 prefix。因此，在 mixed-length trace 上它的 useful rate 不最高，但 recovered saved 更高。

\begin{figure}[H]
\centering
\MaybeFigure[width=0.88\textwidth]{offload_online/04_decision_audit.png}
\caption{CARGO replay 的决策审计：GPU hit、CPU reload 与 cold miss 的组成。该图对应真实输出 \texttt{figures/offload\_online/04\_decision\_audit.png}。}
\label{fig:decision-audit}
\end{figure}

\section{LMCache 真实系统探索}

策略 replay 只能证明 offload 决策在 trace 层面有价值，不能替代真实系统实验。因此我也尝试了 LMCache MP mode。由于当前 vLLM、torch、triton 和 CUDA 版本已经验证可用，我没有直接在原始 \texttt{.venv} 中裸装 LMCache，而是先冻结环境并用 constraints 防止 resolver 自动升级或降级关键包。实际部署时，LMCache server 单独启动，vLLM 通过 \texttt{LMCacheMPConnector} 连接，并在 vLLM 0.20.0 下指定 LMCache-shipped connector module path。

这个路径的正面结果是功能上跑通了：two-request shared-prefix smoke 能在 LMCache 日志中看到 store/retrieve，CoQA pressure 中也能观察到 L1 prefix hit。表 \ref{tab:lmcache} 展示了真实 LMCache baseline 的阶段性结果。g16 下 revisit TTFT 为 38.8ms，接近 GPU warm cache；但 g24/g32 下 revisit TTFT 仍为 350.0ms 和 338.1ms，明显高于 GPU warm，也没有达到理想 CPU offload 性能。

\begin{table}[H]
\centering
\small
\caption{LMCache MP mode 真实系统探索结果}
\label{tab:lmcache}
\begin{tabular}{lcccc}
\toprule
实验文件 & 请求完成数 & Seed TTFT(ms) & Filler TTFT(ms) & Revisit TTFT(ms) \\
\midrule
\texttt{coqa\_pressure\_g16.csv} & 96/96  & 385.9 & 399.5 & 38.8 \\
\texttt{coqa\_pressure\_g24.csv} & 144/144 & 413.9 & 381.9 & 350.0 \\
\texttt{coqa\_pressure\_g32.csv} & 192/192 & 394.2 & 385.8 & 338.1 \\
\bottomrule
\end{tabular}
\end{table}

我认为这部分结果应当如实解释，而不是强行写成性能优化成功。WSL2 环境下 LMCache 的 GPU transfer 路径出现 CUDA IPC \texttt{invalid resource handle}，实际只能退回 non-GPU transfer。这样虽然证明了 KV 可以跨 LMCache server 和 vLLM 连接器进行存取，但 CPU reload 的路径开销太高，抵消了 prefix prefill saved。g24/g32 的 revisit TTFT 相比 seed/filler 有一定下降，但距离 GPU Prefix Cache warm baseline 的 36ms 量级仍然很远。

因此，LMCache 在本报告中的定位是“真实系统功能验证和风险诊断”。它说明后续如果要把 replay 策略变成真实系统收益，必须在 native Linux 或更稳定的 CUDA IPC 环境中继续验证，并进一步解决 CPU/GPU KV 传输、异步 prefetch 和调度时机问题。

\section{结论与后续工作}

本次实验得到的第一个结论是，vLLM 原生 Prefix Cache 在真实服务模式下确实有效，但收益依赖于长前缀和真实复用。短前缀或 cold 请求不应被用来证明 Prefix Cache 有收益；真正稳定的收益出现在长 prefix 的 warm 请求上。CoQA 1536-token 前缀中，warm median saved 达到 38.37ms，说明在文档问答场景中复用 KV 是有实际意义的。

第二个结论是，GPU-only Prefix Cache 存在明显容量 cliff。g16 到 g24 的变化使 revisit useful-hit rate 从 1.0 降到 0，这比平均延迟变化更能说明问题本质：GPU cache 容量不足时，系统并不是逐渐变差，而是会在工作集超过阈值后把大量有价值前缀驱逐掉。这个 cliff 是引入 offload 的直接动机。

第三个结论是，offload 策略需要 utility-aware，而不能只看命中率。Mixed-length trace 上 size-aware 的 useful-hit rate 高于 CARGO，但 recovered saved 远低于 CARGO，说明保留更多“小而容易命中”的 prefix 并不一定带来更大延迟收益。CARGO 通过公式 \ref{eq:cargo-score} 同时考虑 saved、bytes、历史热度和 reload 成本，因此更符合“有限 GPU/CPU 空间应服务于最大 TTFT 净收益”的目标。

本阶段仍有明显限制。Replay 以 prefix group 为粒度，尚未修改 vLLM 内部 block manager 或实现真实 KV tensor reload；CPU reload 成本也只是用固定带宽模型估计。LMCache 虽然功能跑通，但受 WSL2 CUDA IPC 限制，真实性能结果不能作为最终 offload 加速结论。后续工作应当优先在 native Linux 环境中复现 LMCache GPU transfer，并把 CARGO 的 online score 接入真实 cache admission/eviction 路径；在此基础上，再考虑 queue-aware prefetch、时间衰减热度和跨会话 workload。

\section{AI 工具使用说明}

本项目中我使用 AI 工具辅助完成代码组织、实验脚本整理、图表生成逻辑检查和报告初稿润色。实验设计、关键参数选择、vLLM/LMCache 环境调试、CoQA workload 构造、正式实验运行以及结果判断均由我在本地环境中完成和确认。对于 AI 生成或辅助修改的代码，我通过静态检查、smoke run、正式实验日志和图表结果进行验证；对于报告中的实验结论，我只采用本项目实际产生的 CSV、JSONL 和图表，不使用未跑出的假设数据。

\begin{thebibliography}{9}
\bibitem{vllm}
Kwon, W. et al. Efficient Memory Management for Large Language Model Serving with PagedAttention. SOSP 2023.
\bibitem{coqa}
Reddy, S., Chen, D., Manning, C. D. CoQA: A Conversational Question Answering Challenge. TACL 2019.
\bibitem{lmcache}
LMCache Project. LMCache documentation and vLLM connector quickstart.
\bibitem{qwen}
Qwen Team. Qwen2.5 Technical Report and model documentation.
\end{thebibliography}

\end{document}
